\documentclass[12pt]{article}
\usepackage[utf8]{inputenc}
\usepackage[T1]{fontenc}
\usepackage{hyperref}
\title{Tause Language Resources for K-12 Curriculum Development}
\author{Compiled by Claude Sonnet 5 (claude-sonnet-5) for Andr\'as Kornai}
\date{}
\begin{document}
\maketitle

\section{Introduction}
Tause (ISO 639-3 \texttt{tad}, Glottocode taus1252) is a Lakes Plain language of Papua, Indonesia. It is also called Doa or Darha. Sources give about 350 to 500 speakers, mostly in Derapos village, and most speakers are monolingual. Outside contact began only in 1982. The language is poorly documented, and the search found no educational materials.

\section{Primary Educational Resources}
None found.

\section{Linguistic and Reference Materials}
\begin{itemize}
\item Clouse (1997), the main lexical source, with Tause word lists in comparative Lakes Plain tables: \url{https://transnewguinea.org/source/clouse-1997?page=9}
\item Wikipedia overview, covering classification and dialects: \url{https://en.wikipedia.org/wiki/Tause_language}
\item Rosetta Project entry: \url{https://rosettapanglossia.longnow.org/index.php?title=Tad}
\end{itemize}

\section{Religious and Community Resources}
Joshua Project reports printed Bible portions only, and no gospel recordings or films in Tause: \url{https://joshuaproject.net/people_groups/15293}

\section{Digital and Audio Resources}
None found.

\section{Curriculum Development Context}
Almost no materials exist, so any curriculum would depend on new documentation work with the community.
\end{document}